\section{The W state}\label{sec:w_state}
%% ===================================================================

The W state on $N$ sites is the single-excitation symmetric state
\begin{align}
    \ket{W_N}
     & = \ket{10\cdots0} + \ket{010\cdots0} + \cdots + \ket{0\cdots01},
    \notag
\end{align}
the equal-weight superposition of all configurations with exactly one
excitation~\cite{DurVidalCirac2000}. Following
\cite[Appendix~A, ``The W state'']{Cirac2021Matrix}, it is presented with open
boundary conditions, where the matrix product is closed by a left and a right
boundary vector rather than by a trace. The tensor of
Definition~\ref{def:w_tensor} is the one printed there, in the same basis and
with the same boundary vectors, so the results below are stated for it
directly.

\begin{definition}[Open-boundary contraction]\label{def:open_coeff}
    \lean{MPSTensor.openCoeff}
    \leanok
    \uses{def:mps_tensor, def:eval_word}
    For a tensor $A$ with bond dimension $D$, a left boundary covector
    $v_L\in\C^D$, a right boundary vector $v_R\in\C^D$, and a word
    $w=(i_1,\ldots,i_k)$, the open-boundary contraction is the bilinear pairing
    \begin{align}
        (v_L| A^{i_1}A^{i_2}\cdots A^{i_k}\,|v_R)
         & = v_L^{\mathsf T}(A^{i_1}\cdots A^{i_k})v_R.
        \notag
    \end{align}
\end{definition}

\begin{definition}[Open-boundary state]\label{def:open_state}
    \lean{MPSTensor.openState}
    \leanok
    \uses{def:open_coeff}
    The open-boundary state on $N$ sites assigns to each configuration
    $\sigma=(\sigma_1,\ldots,\sigma_N)$ the amplitude
    $(v_L| A^{\sigma_1}\cdots A^{\sigma_N}\,|v_R)$.
\end{definition}

\begin{definition}[W tensor and boundary vectors]\label{def:w_tensor}
    \lean{MPSTensor.wTensor}
    \lean{MPSTensor.wLeftBoundary}
    \lean{MPSTensor.wRightBoundary}
    \lean{MPSTensor.wRaising}
    \leanok
    \uses{def:mps_tensor}
    Following \cite[Appendix~A, ``The W state'']{Cirac2021Matrix}, set $d=D=2$.
    The W tensor is the site-independent family
    \begin{align}
        A^0 & = \begin{pmatrix}1&0\\0&1\end{pmatrix},
        \notag                                        \\
        A^1 & = \begin{pmatrix}0&1\\0&0\end{pmatrix},
        \notag
    \end{align}
    with left boundary covector $(l|=(0|$ and right boundary vector $|r)=|1)$.
    Here $A^1$ is the single raising operator: it sends $|1)\mapsto|0)$ and
    annihilates $|0)$, and squares to zero.
\end{definition}

\begin{definition}[W-state indicator]\label{def:w_indicator}
    \lean{MPSTensor.wIndicator}
    \leanok
    For $N\geq 0$, the W-state amplitude is the indicator of the
    single-excitation configurations:
    \begin{align}
        \ket{W_N}(\sigma_1,\ldots,\sigma_N)
         & =
        \begin{cases}
            1, & \text{exactly one } \sigma_j=1, \\
            0, & \text{otherwise}.
        \end{cases}
        \notag
    \end{align}
\end{definition}

\begin{theorem}[The open-boundary contraction is the W state]%
    \label{thm:w_state_identification}
    \lean{MPSTensor.wTensor_openState_eq_wIndicator}
    \leanok
    \uses{def:w_tensor, def:open_state, def:w_indicator}
    For every $N$, the open-boundary state of the W tensor with boundary
    vectors $(0|$ and $|1)$ equals the W state:
    \begin{align}
        \ket{W_N}
         & = \sum_{\sigma}(0| A^{\sigma_1}\cdots A^{\sigma_N}\,|1)
        \ket{\sigma}.
        \notag
    \end{align}
\end{theorem}

\begin{proof}
    \leanok
    \uses{def:w_tensor, def:open_state}
    Because $A^0$ is the identity and $A^1$ is the raising operator with
    $A^1A^1=0$, the ordered product
    $A^{\sigma_1}\cdots A^{\sigma_N}$ depends only on the number of excitations
    of $\sigma$: it is the identity with no excitation, the raising operator
    with one excitation, and zero with two or more. Reading off the $(0,1)$
    matrix element against the boundary vectors therefore gives $1$ on the
    single-excitation configurations and $0$ elsewhere, which is the W-state
    amplitude.
\end{proof}

\begin{theorem}[Three-site single excitations]%
    \label{thm:w_state_three_excited}
    \lean{MPSTensor.wState_three_excitedAt}
    \lean{MPSTensor.wState_three_vacuum}
    \lean{MPSTensor.wState_three_double}
    \leanok
    \uses{thm:w_state_identification}
    On three sites, each of the three single-excitation configurations
    $\ket{100}$, $\ket{010}$, $\ket{001}$ receives W-state amplitude $1$, while
    the vacuum $\ket{000}$ and the doubly-excited $\ket{110}$ receive
    amplitude $0$.
\end{theorem}

\begin{proof}
    \leanok
    \uses{thm:w_state_identification}
    Direct instances of Theorem~\ref{thm:w_state_identification}: the excitation
    count is one in the first three cases and zero or two in the others.
\end{proof}

The representation of Definition~\ref{def:w_tensor} is not translationally
invariant, because its boundary is not periodic
\cite[Appendix~A, ``The W state'']{Cirac2021Matrix}. The following computation
illustrates this remark by closing the same tensor with a trace.

\begin{theorem}[The periodic closure is not the W state]%
    \label{thm:w_state_periodic_closure}
    \lean{MPSTensor.mpv_wTensor}
    \lean{MPSTensor.mpv_wTensor_eq_zero_of_wIndicator_eq_one}
    \lean{MPSTensor.mpv_wTensor_ne_wIndicator}
    \leanok
    \uses{def:w_tensor, def:mpv, def:w_indicator}
    Closing the tensor of Definition~\ref{def:w_tensor} with a trace gives
    \begin{align}
        \operatorname{tr}\bigl(A^{\sigma_1}\cdots A^{\sigma_N}\bigr)
         & =
        \begin{cases}
            2, & \text{no } \sigma_j=1, \\
            0, & \text{otherwise}.
        \end{cases}
        \notag
    \end{align}
    In particular the trace contraction vanishes on every single-excitation
    configuration, and for $N\geq 1$ it is not $\ket{W_N}$.
\end{theorem}

\begin{proof}
    \leanok
    \uses{def:w_tensor, def:mpv, def:w_indicator}
    The ordered product is the identity, the raising operator, or zero according
    to the number of excitations, as in the proof of
    Theorem~\ref{thm:w_state_identification}. Their traces are $2$, $0$, and $0$.
    For $N\geq1$ the configuration $\ket{10\cdots0}$ has W-state amplitude $1$
    and trace contraction $0$.
\end{proof}

\begin{definition}[Periodic representation of the W state]%
    \label{def:w_state_periodic_rep}
    \lean{MPSTensor.IsPeriodicWState}
    \leanok
    \uses{def:mpv, def:w_indicator}
    A tensor $B$ with physical dimension $2$ and bond dimension $D$ is a periodic
    representation of $\ket{W_N}$ when
    $\operatorname{tr}(B^{\sigma_1}\cdots B^{\sigma_N})=\ket{W_N}(\sigma)$ for every
    configuration $\sigma$ of $N$ sites.
\end{definition}

\begin{theorem}[Bond dimension one]%
    \label{thm:w_state_periodic_bond_one}
    \lean{MPSTensor.not_isPeriodicWState_of_bondDim_one}
    \leanok
    \uses{def:w_state_periodic_rep}
    For $N\geq 2$, no tensor of bond dimension $1$ is a periodic representation of
    $\ket{W_N}$.
\end{theorem}

\begin{proof}
    \leanok
    \uses{def:w_state_periodic_rep}
    Write $B^0=b_0$ and $B^1=b_1$ as scalars. The vacuum amplitude gives
    $b_0^N=0$, so $b_0=0$. The configuration $\ket{10\cdots0}$ then has amplitude
    $b_1b_0^{N-1}=0$, whereas its W-state amplitude is $1$.
\end{proof}

\begin{theorem}[No fixed bond dimension represents every W state]%
    \label{thm:w_state_periodic_bound}
    \lean{MPSTensor.trace_pow_eq_zero_of_isPeriodicWState}
    \lean{MPSTensor.pow_eq_zero_of_isPeriodicWState}
    \lean{MPSTensor.not_isPeriodicWState_of_pow_eq_zero}
    \lean{MPSTensor.not_isPeriodicWState_of_lt}
    \lean{MPSTensor.exists_not_isPeriodicWState_le}
    \lean{MPSTensor.not_forall_isPeriodicWState}
    \leanok
    \uses{def:w_state_periodic_rep}
    Let $B$ be a tensor with physical dimension $2$ and bond dimension $D$. If $B$
    is a periodic representation of $\ket{W_k}$ for every $1\leq k\leq D$, then
    $(B^0)^D=0$, and $B$ is a periodic representation of $\ket{W_N}$ for no
    $N>D$. Consequently there is a length $1\leq N\leq D+1$ for which $B$ is not
    a periodic representation of $\ket{W_N}$, and no single tensor of fixed bond
    dimension is a periodic representation of $\ket{W_N}$ for every $N\geq 1$.
\end{theorem}

\begin{proof}
    \leanok
    \uses{def:w_state_periodic_rep}
    The vacuum amplitudes give $\operatorname{tr}((B^0)^k)=0$ for
    $1\leq k\leq D$. By the Newton--Girard identities the characteristic
    polynomial of $B^0$ is $x^D$, and the Cayley--Hamilton theorem gives
    $(B^0)^D=0$. For $N>D$ the configuration $\ket{10\cdots0}$ then has
    amplitude $\operatorname{tr}(B^1(B^0)^{N-1})=0$, whereas its W-state
    amplitude is $1$. If $B$ represented every $\ket{W_N}$ with
    $1\leq N\leq D+1$, the case $N=D+1$ would contradict this.
\end{proof}

\begin{lemma}[Cut rank of the W state]\label{lem:w_state_cut_rank}
    \lean{MPSTensor.weightIndicator, MPSTensor.wIndicator_append,
        MPSTensor.wIndicator_append_mem_span,
        MPSTensor.finrank_span_wIndicator_append_le_two}
    \leanok
    \uses{def:w_indicator}
    Split $N=R+R'$ sites into a first piece of $R$ sites and a second piece of
    $R'$ sites. For every configuration $\tau$ of the second piece, the vector
    $\sigma\mapsto\braket{\sigma\tau}{W_N}$ on the first piece is a combination of
    the vacuum indicator and the one-excitation indicator $\ket{W_R}$. Hence the
    reduced state of $\ket{W_N}$ on the first piece has rank at most two.
\end{lemma}

\begin{proof}
    \leanok
    If $\tau$ contains one excitation, the amplitude of $\sigma\tau$ is one exactly
    when $\sigma$ is the vacuum; if $\tau$ is the vacuum, it is one exactly when
    $\sigma$ has one excitation; otherwise it vanishes.
\end{proof}

\begin{theorem}[Single-length bound under a canonical form with condition C1]%
    \label{thm:w_state_canonical_bound}
    \lean{MPSTensor.PGVWC07CanonicalFormData.lt_of_mpv_eq_smul_wIndicator,
        MPSTensor.PGVWC07CanonicalFormData.lt_of_mpv_eq_smul_wIndicator_of_isUnit,
        MPSTensor.PGVWC07CanonicalFormData.lt_of_mpv_eq_smul_wIndicator_of_isPrimitive}
    \leanok
    \uses{def:pgvwc07_ti_canonical_representation, def:l_blk_injective,
        def:blocks_not_gauge_phase_equiv, def:w_indicator}
    Let $A$ be a translation-invariant tensor of bond dimension $D$ in the
    canonical form of Definition~\ref{def:pgvwc07_ti_canonical_representation},
    with $b$ blocks, no two of which are related by a gauge transformation and a
    phase. Suppose that for some $N$ and some $c\neq0$ the periodic vector of $A$
    on $N$ sites is $c\ket{W_N}$.
    \begin{enumerate}
        \item If every block is $L_0$-block injective for a common $L_0>0$, then
              $N<2\max\bigl(L_0,\,3(b-1)(L_0+1)\bigr)$.
        \item If every block is $L$-block injective for some $L>0$ and its first
              matrix $A^0$ is invertible, then
              $N<2\max\bigl(D^2+1,\,3(D-1)(D^2+2)\bigr)$, so $D=\Omega(N^{1/3})$.
        \item If the dual transfer channel of every block is primitive, then
              $N<2\max\bigl((D^2+1)^2,\,3(D-1)((D^2+1)^2+1)\bigr)$.
    \end{enumerate}
    The first item is the inequality of the proof in
    \cite{PerezGarcia2007Matrix}, conditional on the canonical form with condition
    C1, and with the threshold enlarged by $L_0$ so that it also covers a single
    block. The distinctness hypothesis is stricter than the source's
    pairwise-different block states. The reduction of an arbitrary tensor to this
    form and the review's $D^3\log D=\Omega(N)$ are not formalized
    \cite{gap:rmp_w_state_ti_bound}.
\end{theorem}

\begin{proof}
    \leanok
    \uses{lem:w_state_cut_rank}
    Cut the chain into two pieces of at least
    $n_0=\max(L_0,3(b-1)(L_0+1))$ sites. By condition C1 and the direct-sum
    lemma, the reduced state of the first piece has rank at least $\sum_j D_j^2$,
    while Lemma~\ref{lem:w_state_cut_rank} bounds it by two. Hence there are at
    most two blocks, each of size one by one. A sum of at most two
    translation-invariant product states is not a nonzero multiple of
    $\ket{W_N}$ for $N\geq3$, a contradiction when $N\geq2n_0$. For the second
    item, an invertible letter gives condition C1 at $L_0=D^2$ by the fixed-length
    Wielandt theorem; for the third, the general quantum Wielandt bound gives
    $L_0\leq(D^2+1)^2$.
\end{proof}

\begin{remark}[Growth of the translation-invariant bond dimension]%
    \label{rem:w_state_ti_bound}
    \uses{def:w_state_periodic_rep}
    Theorem~\ref{thm:w_state_periodic_bound} concerns one tensor representing
    the W state on a range of lengths. The review states the stronger
    single-length bound: a translation-invariant representation of $\ket{W_N}$
    with bond dimension $D$ requires $D^3\log D=\Omega(N)$, hence
    $D=\Omega(N^{1/(3+\delta)})$ for every $\delta>0$, obtained by combining the
    representation theory of translation-invariant matrix product states with a
    quantitative quantum Wielandt inequality
    \cite[Appendix~A, ``The W state'']{Cirac2021Matrix}. This single-length bound
    is not formalized; Theorem~\ref{thm:w_state_canonical_bound} proves the
    conditional single-length bound of its source~\cite{gap:rmp_w_state_ti_bound}.
\end{remark}

%% ===================================================================
